\documentclass[11pt,a4paper]{ctexart}
\usepackage[margin=2cm]{geometry}
\usepackage{xcolor}
\usepackage{booktabs}
\usepackage{tabularx}
\usepackage{enumitem}
\usepackage{listings}
\usepackage{hyperref}
\usepackage{fancyhdr}

\definecolor{codebg}{HTML}{F4F6F8}
\definecolor{accent}{HTML}{1769AA}
\hypersetup{colorlinks=true,linkcolor=accent,urlcolor=accent}
\setlist{nosep,leftmargin=2em}
\setlength{\parindent}{0pt}
\setlength{\parskip}{0.45em}
\pagestyle{fancy}
\fancyhf{}
\lhead{Linux、Git 与 HTTP 基础}
\rhead{课堂讲义}
\cfoot{\thepage}
\lstset{
  basicstyle=\ttfamily\small,
  backgroundcolor=\color{codebg},
  frame=single,
  rulecolor=\color{codebg},
  breaklines=true,
  columns=fullflexible,
  keepspaces=true,
  showstringspaces=false
}

\newcommand{\key}[1]{\textcolor{accent}{\textbf{#1}}}

\title{\textbf{Linux 基础、Git 与请求状态码}\\\large 一节课入门讲义}
\author{}
\date{}

\begin{document}
\maketitle
\vspace{-2em}

\begin{center}
\fcolorbox{accent}{codebg}{\parbox{0.88\textwidth}{
\textbf{学习目标：}能在 Linux 中定位和操作文件；能用 Git 完成一次基本版本提交；能根据 HTTP 状态码快速判断请求结果。
}}
\end{center}

\section{Linux 基础}

Linux 是一种广泛用于服务器和开发环境的操作系统。命令行的基本形式是：

\begin{lstlisting}
command [options] [arguments]
ls -la /home
\end{lstlisting}

其中 \texttt{ls} 是命令，\texttt{-la} 是选项，\texttt{/home} 是操作对象。

\subsection{目录与路径}

\begin{itemize}
  \item \key{绝对路径}：从根目录 \texttt{/} 开始，例如 \texttt{/home/student/project}。
  \item \key{相对路径}：从当前目录开始，例如 \texttt{./src}。
  \item \texttt{.} 表示当前目录，\texttt{..} 表示上一级目录。
  \item Linux 文件名区分大小写：\texttt{Readme.md} 与 \texttt{README.md} 是两个文件。
\end{itemize}

\subsection{常用命令速查}

\begin{tabularx}{\textwidth}{>{\ttfamily}l X l}
\toprule
\normalfont 命令 & 作用 & 示例 \\
\midrule
pwd & 显示当前目录 & \texttt{pwd} \\
ls & 查看目录内容 & \texttt{ls -la} \\
cd & 切换目录 & \texttt{cd project} \\
mkdir & 创建目录 & \texttt{mkdir demo} \\
touch & 创建空文件或更新时间 & \texttt{touch notes.txt} \\
cp & 复制文件或目录 & \texttt{cp a.txt b.txt} \\
mv & 移动或重命名 & \texttt{mv old.txt new.txt} \\
rm & 删除文件；删除通常难以恢复 & \texttt{rm notes.txt} \\
cat & 输出文件内容 & \texttt{cat notes.txt} \\
grep & 在文本中搜索 & \texttt{grep error app.log} \\
man & 查看命令手册 & \texttt{man ls} \\
\bottomrule
\end{tabularx}

\subsection{权限概念}

\texttt{ls -l} 可能显示 \texttt{-rwxr-xr--}。前三位是所有者权限，中间三位是用户组权限，最后三位是其他人权限。\texttt{r/w/x} 分别表示读、写、执行。

\begin{lstlisting}
chmod +x deploy.sh     # add execute permission
sudo command           # run as administrator (use carefully)
\end{lstlisting}

\textbf{小提醒：}执行命令前确认自己所在的目录；使用 \texttt{rm}、\texttt{sudo} 前尤其要检查目标。

\section{Git：让修改可追踪}

Git 是\key{分布式版本控制系统}。它记录文件的变化，便于回退、协作和审查。

\subsection{三个区域}

\begin{center}
工作区 $\xrightarrow{\texttt{git add}}$ 暂存区 $\xrightarrow{\texttt{git commit}}$ 本地仓库
\end{center}

\begin{itemize}
  \item \key{工作区}：正在编辑的文件。
  \item \key{暂存区}：准备放进下一次提交的修改。
  \item \key{本地仓库}：已经提交的版本历史。
\end{itemize}

\subsection{一次完整的基本流程}

\begin{lstlisting}
git init                         # initialize a repository
git status                       # inspect current state
git add README.md                # stage a file
git commit -m "docs: add intro"  # create a commit
git log --oneline                # view compact history
\end{lstlisting}

提交信息应简短说明“这次修改做了什么”，一次提交尽量只解决一个主题。

\subsection{分支与远程协作}

\begin{lstlisting}
git switch -c feature/login  # create and switch branch
git switch main              # switch back to main
git pull                     # fetch and merge updates
git push                     # push local commits
\end{lstlisting}

\textbf{推荐习惯：}开始工作前先同步；在功能分支上修改；提交前运行 \texttt{git status} 和测试；发生冲突时理解双方修改后再解决。

\section{HTTP 请求状态码}

浏览器或客户端向服务器发送请求，服务器返回响应。状态码是三位数字，用来概括请求结果。\key{第一位数字表示大类}。

\begin{tabularx}{\textwidth}{c l X}
\toprule
类别 & 含义 & 常见判断 \\
\midrule
1xx & 信息性响应 & 请求已收到，仍在处理 \\
2xx & 成功 & 请求已被正常处理 \\
3xx & 重定向 & 资源位置变化，或可使用缓存 \\
4xx & 客户端错误 & 请求参数、身份或权限等有问题 \\
5xx & 服务器错误 & 服务端处理请求时发生异常 \\
\bottomrule
\end{tabularx}

\subsection{常见状态码}

\begin{tabularx}{\textwidth}{c l X}
\toprule
状态码 & 名称 & 典型场景 \\
\midrule
200 & OK & 请求成功并返回内容 \\
201 & Created & 创建资源成功，例如新增用户 \\
204 & No Content & 成功，但没有响应正文 \\
301 & Moved Permanently & 资源永久迁移 \\
302 & Found & 临时重定向 \\
304 & Not Modified & 缓存仍有效，可继续使用 \\
400 & Bad Request & 参数格式或请求语法有误 \\
401 & Unauthorized & 未认证或凭证无效，通常需要登录 \\
403 & Forbidden & 已识别身份，但没有访问权限 \\
404 & Not Found & 请求的资源不存在 \\
409 & Conflict & 当前状态冲突，例如重复创建 \\
429 & Too Many Requests & 请求过于频繁，触发限流 \\
500 & Internal Server Error & 服务器内部异常 \\
502 & Bad Gateway & 网关收到上游服务的无效响应 \\
503 & Service Unavailable & 服务暂时不可用或过载 \\
504 & Gateway Timeout & 网关等待上游服务超时 \\
\bottomrule
\end{tabularx}

\textbf{重点辨析：}\texttt{401} 强调“你是谁尚未通过验证”；\texttt{403} 强调“知道你是谁，但你无权访问”。\texttt{404} 不等于服务器宕机，而是目标资源未找到。

\subsection{排查请求问题的顺序}

\begin{enumerate}
  \item 先看状态码的大类：成功、客户端问题还是服务端问题。
  \item 再看响应正文和响应头中的错误信息。
  \item 检查 URL、请求方法、参数、登录状态与权限。
  \item 若为 5xx，结合服务端日志定位异常或上游依赖。
\end{enumerate}

\section{随堂练习}

\begin{enumerate}
  \item 在终端创建目录 \texttt{lesson-demo}，进入目录并创建 \texttt{README.md}。
  \item 初始化 Git 仓库，把 \texttt{README.md} 加入暂存区并提交。
  \item 某接口返回 401、403、404 时，分别应优先检查什么？
  \item 用户提交表单后成功创建记录，最合适的状态码是什么？
  \item 网关等待后端服务超时，最可能看到哪个状态码？
\end{enumerate}

\textbf{参考答案：}（1--2）依次使用 \texttt{mkdir}、\texttt{cd}、\texttt{touch}、\texttt{git init}、\texttt{git add}、\texttt{git commit}；（3）认证信息、访问权限、资源路径；（4）201；（5）504。

\section*{一页记忆}

\begin{itemize}
  \item Linux：先用 \texttt{pwd} 定位，再用 \texttt{ls} 查看；危险命令先核对路径。
  \item Git：\texttt{status} 看状态，\texttt{add} 选修改，\texttt{commit} 存版本。
  \item HTTP：2xx 成功，3xx 重定向，4xx 查请求，5xx 查服务端。
\end{itemize}

\end{document}
